\section{Stability-Guided Observational Prioritization}

The second stage converts prompt-side variables into a bounded set of
intervention candidates. It combines type-derived exclusions with FCI evidence
that preserves endpoint uncertainty, then tests whether candidate relations
remain stable when whole semantic task clusters are resampled. FCI is one
budgeted selector; held-out randomization, not a PAG endpoint, identifies the
prompt-policy effect.

\subsection{TSG-Derived Constraints}

The Prompt TSG supplies variables, background knowledge, and edge constraints;
it does not supply learned causal edges. Temporal tiers prohibit outcomes from
causing pre-treatment metadata or prompt factors. Catalog types define
family- and scope-specific forbidden directions, and typed adjacency
restrictions exclude semantically impossible local relations. These rules may
forbid both directions of an adjacency, but they never require a
feature--outcome edge or any edge on the possible path being tested. \model
persists the exact tiers, forbidden pairs, typed adjacency matrix, variable
order, and background-knowledge digest and validates every learned PAG against
them.

\subsection{FCI and Semantic-Cluster Stability}

The minimum backend is the established causal-learn implementation of
FCI~\cite{spirtes1995fci,zheng2024causallearn}. Its PAG output preserves
endpoint uncertainty under latent-variable assumptions rather than selecting a
single convenient orientation~\cite{zhang2008ancestral}. The run locks the
G-square conditional-independence test, variable order, significance level,
depth and path bounds, backend version, and background knowledge.

The fixed-reference analysis selects one frozen request-randomness slot per
task. The two-level analysis resamples \texttt{semantic\_task\_cluster\_id}
and then selects one frozen request slot inside each sampled task occurrence.
A separately frozen multi-slot sensitivity aggregates slots at task level
without passing fractional pseudo-categories to G-square. None treats repeated
generations, task instances, or rows as independent clusters.

Candidate extraction retains PAG endpoint marks and scores permitted
feature--outcome adjacency, endpoint-compatible possible ancestry, or the same
relation within the frozen $C_q$-eligible population. Simultaneously authored
Prompt features occupy one temporal tier: their within-tier endpoints are not
interpreted as security mechanisms or mediation. \model applies a fixed
top-$K$ budget and deterministic ordering rather than claiming recovery of a
unique DAG.

\subsection{Hypothesis Freeze}

Before any selector runs, one \texttt{CandidateUniverseManifest} freezes
skeletons $\widetilde h=(C_q,f,a,\mathcal Q_h,Y)$ and their common information
budget. After all rankings are frozen, the unique top-$K$ union crosses one
selector-invariant, outcome-blind bridge. A successful bridge produces
$h=(C_q,f,a,Q_h,Y)$, binding one non-actionable context, one actionable
feature, one atomic operation, a finite realization distribution, and one
outcome. Its record also binds CWE/model scope, reference PAG endpoints,
semantic-cluster support, expected direction, and the table, catalog,
extractor, selector, CI-test, and background-knowledge provenance.

This hypothesis freeze is committed before any confirmation variant,
assignment, generated code, or outcome exists. Held-out evidence may classify
a frozen hypothesis, but it cannot add, remove, rerank, remap, or rewrite one.
An empty, bridge-failed, or protocolization-failed selector slot remains a zero
in the original budget denominator.
